\documentclass[11pt]{article}
\usepackage[margin=1in]{geometry}
\usepackage{amsmath,amssymb}
\usepackage{graphicx}
\usepackage{booktabs}
\usepackage{hyperref}
\usepackage{xcolor}

\title{M1M3 Mirror Thermal Response to Setpoint Changes:\\
Surface-Resolved Analysis}
\author{J.\ Esteves}
\date{\today}

\begin{document}
\maketitle

%----------------------------------------------------------------------
\section{Introduction}

The M1M3 thermal control system maintains the mirror at a target
temperature below ambient to prevent convective ``mirror seeing.''
The Environmental Awareness Service (EAS) issues setpoint commands to
the glycol heater system (\texttt{MTM1M3TS.command\_applySetpoints},
field \texttt{heatersSetpoint}), which heats or cools the mirror
through glycol circulation.

This report characterizes the actual thermal transfer function from
setpoint to mirror temperature using 149 observing nights of archival
data (July 2025 -- February 2026).  The goals are:
\begin{enumerate}
  \item Document the EAS setpoint strategy (offset, rate, cadence).
  \item Measure the glycol system transfer function (lag, time constant).
  \item Validate and calibrate our first-order thermal model
    $dT/dt = (T_\mathrm{sp} - T_\mathrm{mirror})/\tau$.
  \item Quantify the actual EAS tracking performance.
  \item Resolve the thermal behavior at three depth levels---back
    (near FCUs), middle, and front (optical surface)---using the
    thermocouple gradient-fit decomposition.
  \item Characterize the coupling of each surface with cell air
    (ESS~113) and glycol, in both pre-control and active-control
    periods.
\end{enumerate}

%----------------------------------------------------------------------
\section{Data}

All data come from the pre-built observing-night parquet file
(166~nights, 1-min resolution) and the thermocouple gradient dataset
(164~nights).  Key channels:

\begin{table}[htbp]
\centering
\caption{Sensor channels used in this analysis.}
\label{tab:channels}
\begin{tabular}{llc}
\toprule
Channel & Source & Coverage \\
\midrule
\texttt{heatersSetpoint\_mean} & \texttt{MTM1M3TS.command\_applySetpoints} & 8.9\% (163 nights) \\
$T_\mathrm{mid}$ (intercept + 0.5$\times$z\_gradient) & 96-thermocouple gradient fit & 97\% (164 nights) \\
\texttt{z\_gradient}, \texttt{z\_gradient\_err} & 96-thermocouple gradient fit & 97\% \\
\texttt{mirrorCoolantSupply/Return} & \texttt{MTM1M3TS.glycolLoopTemperature} & 97\% \\
\texttt{inside\_m1m3\_temp\_mean} & ESS~113 & 99.5\% \\
\texttt{outside\_temp\_mean} & ESS~301 & 99.7\% \\
\bottomrule
\end{tabular}
\end{table}

The \textbf{mirror temperature reference} is the mid-surface temperature
$T_\mathrm{mid}$ ($z_\mathrm{norm} = 0.5$) from the gradient-fit
decomposition of the 96 thermocouples (see Section~\ref{sec:sensor_map}).
This is the best proxy for the bulk glass temperature, as it lies
halfway between the glycol-cooled back surface and the air-facing
optical surface.  The gradient \texttt{intercept} (back surface center,
$z_\mathrm{norm} = 0$) and \texttt{z\_gradient\_err} (through-glass
gradient uncertainty) are also used.

The setpoint is a sparse signal (issued every $\sim$3\,min as discrete
commands); it is forward-filled onto the 1-min grid for time-series
analysis.  Nights with fewer than 10 setpoint commands or fewer than
60~min of mirror temperature data are excluded, leaving 149 valid nights.

%----------------------------------------------------------------------
\section{EAS Setpoint Strategy}

\subsection{Setpoint Offset}

The EAS targets the heater setpoint at a fixed offset below the M1M3
air temperature (ESS~113):

\begin{table}[htbp]
\centering
\caption{Setpoint offset statistics (149 nights).}
\label{tab:offset}
\begin{tabular}{lccc}
\toprule
Offset & Median & Mean & Std \\
\midrule
Setpoint $-$ ESS~113 & $-1.00$\,\textdegree C & $-1.36$\,\textdegree C & 1.27\,\textdegree C \\
Setpoint $-$ $T_\mathrm{mid}$ ($z_\mathrm{norm}=0.5$) & $-0.07$\,\textdegree C & $-0.35$\,\textdegree C & 1.15\,\textdegree C \\
\bottomrule
\end{tabular}
\end{table}

The EAS uses a \textbf{1.0\,\textdegree C cold bias} relative to
ESS~113, not the 0.3\,\textdegree C assumed in our simulation model.
The setpoint sits only $-0.07$\,\textdegree C below $T_\mathrm{mid}$,
meaning the glycol system maintains the bulk glass temperature almost
exactly at the commanded setpoint.
Figure~\ref{fig:strategy} (top left) shows the offset distributions.

\subsection{Setpoint Rate}

The EAS changes the setpoint at a median rate of 0.6\,\textdegree C/h,
with a 95th percentile of 1.8\,\textdegree C/h
(Figure~\ref{fig:strategy}, top right).  Notably, the EAS
\textbf{does not respect the 0.7\,\textdegree C/h rate limit} that we
enforce in our simulation---roughly half of the EAS rate changes exceed
this threshold.  This suggests the hardware can sustain higher rates
than we assumed.

\subsection{Update Cadence}

Setpoint commands are issued every 3.0\,min (median), with occasional
gaps of 10--15\,min (Figure~\ref{fig:strategy}, bottom left).  This is
much faster than our 15-min simulation timestep.

\begin{figure}[htbp]
\centering
\includegraphics[width=\textwidth]{../figures/mirror_setpoint_strategy.png}
\caption{EAS setpoint strategy.  Top left: offset distributions
  (setpoint minus ESS~113 and minus $T_\mathrm{mid}$).  Top right:
  absolute setpoint rate distribution; the red line marks our assumed
  0.7\,\textdegree C/h limit.  Bottom left: update cadence (median
  3\,min).  Bottom right: example night showing all thermal channels
  including $T_\mathrm{mid}$ ($z_\mathrm{norm} = 0.5$).}
\label{fig:strategy}
\end{figure}

%----------------------------------------------------------------------
\section{Glycol System Transfer Function}

\subsection{Heat Flux Proxy}

The coolant supply--return temperature difference
($T_\mathrm{supply} - T_\mathrm{return}$) is a proxy for heat flux
through the glycol loop.  Figure~\ref{fig:transfer} (top left) shows
that when the setpoint is far below the mirror temperature (large
negative error), the supply--return $\Delta T$ reaches $-6$ to
$-7$\,\textdegree C, indicating strong active cooling.  Near zero
setpoint error, the supply--return $\Delta T$ clusters near zero
(passive tracking).

\subsection{Setpoint-to-Mirror Lag}

Cross-correlating the setpoint rate with the $T_\mathrm{mid}$
rate gives a median lag of \textbf{4\,minutes}
(Figure~\ref{fig:transfer}, top right).  This is the time for a
setpoint change to begin affecting the bulk glass temperature.  The
short lag confirms that the glycol system responds rapidly to setpoint
commands.

\subsection{Setpoint Tracking}

Figure~\ref{fig:transfer} (bottom right) confirms that the EAS tracks
ESS~113 with a $-1$\,\textdegree C offset: the per-night trajectories
cluster tightly along the $\mathrm{ESS\,113} - 1$\,\textdegree C line.

\begin{figure}[htbp]
\centering
\includegraphics[width=\textwidth]{../figures/mirror_step_response.png}
\caption{Glycol transfer function.  Top left: supply--return $\Delta T$
  vs setpoint error relative to $T_\mathrm{mid}$ (hexbin density).  Top
  right: cross-correlation lag between setpoint rate and $T_\mathrm{mid}$ rate.
  Bottom left: mirror response to setpoint step events.  Bottom right:
  per-night setpoint vs ESS~113 scatter (20 nights shown).}
\label{fig:transfer}
\end{figure}

%----------------------------------------------------------------------
\section{Model Validation}

\subsection{First-Order Model}

We test the standard first-order thermal model:
\begin{equation}
  \frac{dT_\mathrm{mirror}}{dt}
    = \frac{T_\mathrm{setpoint} - T_\mathrm{mirror}}{\tau}
\end{equation}
by forward-simulating from the actual EAS setpoint trajectory and
comparing with $T_\mathrm{mid}$ ($z_\mathrm{norm} = 0.5$).  The
simulation uses 1-min timesteps with the actual (forward-filled)
setpoint as input.

\subsection{Time Constant Estimation}

We sweep $\tau$ from 0.25 to 6.0\,h in 0.25\,h steps and compute the
RMSE against the mirror mean for each night.

\begin{table}[htbp]
\centering
\caption{Time constant estimation results (149 nights).}
\label{tab:tau}
\begin{tabular}{lc}
\toprule
Metric & Value \\
\midrule
Per-night median $\tau_\mathrm{glycol}$ (vs $T_\mathrm{mid}$) & 2.50\,h \\
Per-night median $\tau_\mathrm{air}$ & 12.0\,h \\
Median RMSE glycol-only & 0.059\,\textdegree C \\
\bottomrule
\end{tabular}
\end{table}

Using $T_\mathrm{mid}$ as the reference temperature (instead of the
back surface or the AboveMirror glycol sensor) gives well-constrained
fits:
\begin{itemize}
  \item The glycol-only model achieves median RMSE = 0.059\,\textdegree C
    with $\tau_\mathrm{glycol}$ median = 2.5\,h.
  \item The two-input model (glycol + air convection) yields
    $\tau_\mathrm{glycol}$ median = 2.5\,h and $\tau_\mathrm{air}$
    median = 12\,h.
  \item $T_\mathrm{mid}$ represents the bulk glass temperature at
    $z_\mathrm{norm} = 0.5$, midway between the glycol-cooled back
    surface and the air-facing optical surface.
\end{itemize}

The shorter $\tau_\mathrm{glycol} = 2.5$\,h compared to the back
surface value of 3.0\,h (Table~\ref{tab:tau_surface}) reflects the
additional air-convection coupling experienced by the mid-level glass.
The first-order model captures the bulk dynamics with high fidelity.

\subsection{Recommended Time Constant}

The optimal $\tau$ for the bulk glass ($T_\mathrm{mid}$) is
\textbf{2.5\,h}, with surface-dependent values ranging from
$\tau_\mathrm{back} = 3.0$\,h (glycol-dominated) to
$\tau_\mathrm{front} = 2.0$\,h (air-coupled).  We recommend using
$\tau = 2.5$\,h as the effective bulk time constant in the simulation
model.

\begin{figure}[htbp]
\centering
\includegraphics[width=\textwidth]{../figures/mirror_model_validation.png}
\caption{Model validation using $T_\mathrm{mid}$.  Top left:
  $\tau_\mathrm{glycol}$ vs $\tau_\mathrm{air}$ per night (colored by
  RMSE).  Top right: RMSE comparison between glycol-only and two-input
  models.  Bottom left: example night comparing model fits.
  Bottom right: per-night $\tau$ distributions.}
\label{fig:validation}
\end{figure}

\begin{figure}[htbp]
\centering
\includegraphics[width=\textwidth]{../figures/mirror_thermal_overview.png}
\caption{Four example nights showing setpoint, ESS~113, $T_\mathrm{mid}$
  ($\pm z\_gradient\_err/2$ shaded band), coolant supply/return, and the
  first-order model fit.  Top left: lowest $\tau$ night.  Top right:
  highest $\tau$ night.  Bottom left: best-fit night.
  Bottom right: worst-fit night.}
\label{fig:overview}
\end{figure}

%----------------------------------------------------------------------
\section{EAS Actual Performance}

We evaluate the current EAS controller performance by comparing
$T_\mathrm{mid}$ against the target temperature
$T_\mathrm{target} = \mathrm{ESS\,113} - 1.0$\,\textdegree C.

\begin{table}[htbp]
\centering
\caption{EAS tracking performance (149 nights).}
\label{tab:performance}
\begin{tabular}{lc}
\toprule
Metric & Value \\
\midrule
Tracking error median ($T_\mathrm{mid}$ $-$ target) & $-0.23$\,\textdegree C \\
Tracking error RMS (all timesteps) & 1.10\,\textdegree C \\
Per-night RMS median & 0.74\,\textdegree C \\
Per-night $\pm$0.3\,\textdegree C compliance median & 26.5\% \\
EAS setpoint rate p95 & 1.80\,\textdegree C/h \\
\bottomrule
\end{tabular}
\end{table}

The error distribution (Figure~\ref{fig:performance}, top left) is
asymmetric, skewed toward negative values (mirror colder than target),
consistent with the cold-is-preferred design philosophy.  The per-night
RMS of 0.74\,\textdegree C (Figure~\ref{fig:performance}, bottom left)
reflects the fundamental challenge: the mirror's 2.5\,h thermal time
constant limits how fast it can track ambient changes.

The $\pm$0.3\,\textdegree C compliance of 26.5\% (Figure~\ref{fig:performance},
bottom right) shows significant room for improvement.  The setpoint rate
distribution (Figure~\ref{fig:performance}, top right) confirms that the
EAS already operates beyond the 0.7\,\textdegree C/h limit we impose
in simulation.

\begin{figure}[htbp]
\centering
\includegraphics[width=\textwidth]{../figures/mirror_eas_performance.png}
\caption{EAS actual performance.  Top left: tracking error distribution
  ($T_\mathrm{mid} - (\mathrm{ESS\,113} - 1.0)$\,\textdegree C).  Top right:
  EAS setpoint rate distribution with $\pm$0.7\,\textdegree C/h markers.
  Bottom left: per-night RMS distribution.  Bottom right: per-night
  $\pm$0.3\,\textdegree C compliance.}
\label{fig:performance}
\end{figure}

%----------------------------------------------------------------------
\section{M1M3 Sensor Map and Gradient-Fit Decomposition}
\label{sec:sensor_map}

The M1M3 mirror is instrumented with 146 thermocouples distributed
across four scanners (ESS~114--117), three depth levels (Front, Middle,
Back), and six azimuthal sectors (A--F).  The glycol thermal control
system uses 96~Fan Coil Units (FCUs) on the back surface of the mirror.
ESS~113 is a humidity/temperature sensor measuring the cell air above
the mirror---it is \emph{not} on the glass surface.

Figure~\ref{fig:sensormap} shows the spatial layout.  The left panel
is the top view ($x$--$y$ plane), showing thermocouple locations
colored by depth level and FCU positions.  The right panel is the
side view (radial--$z_\mathrm{norm}$ plane), illustrating the glass
cross-section with the M3 region (thinner, inner) and M1 region
(thicker, outer).

\subsection{Gradient-Fit Model}

At each 1-min timestamp, a least-squares fit to the 96-thermocouple
temperatures yields:
\begin{equation}
  T(x, y, z_\mathrm{norm})
    = \mathrm{intercept}
    + x\_gradient \cdot x
    + y\_gradient \cdot y
    + z\_gradient \cdot z_\mathrm{norm}
  \label{eq:gradfit}
\end{equation}
where $z_\mathrm{norm} = 0$ (back, near FCUs), $0.5$ (middle), and
$1.0$ (front, optical surface).  The intercept is the temperature at
the mirror center at $z_\mathrm{norm} = 0$, i.e., the \textbf{back
surface center temperature}.

We define three virtual surface temperatures:
\begin{align}
  T_\mathrm{back}  &= \mathrm{intercept}
    & (z_\mathrm{norm} = 0) \\
  T_\mathrm{mid}   &= \mathrm{intercept} + 0.5 \times z\_gradient
    & (z_\mathrm{norm} = 0.5) \\
  T_\mathrm{front} &= \mathrm{intercept} + z\_gradient
    & (z_\mathrm{norm} = 1.0)
\end{align}
The $z\_gradient$ represents the front-minus-back temperature
difference through the glass.  The fit also provides
$z\_gradient\_err$, the $1\sigma$ uncertainty on the through-glass
gradient.

\begin{figure}[htbp]
\centering
\includegraphics[width=\textwidth]{../figures/m1m3_sensor_map.png}
\caption{M1M3 sensor layout.  Left: top view showing 146
  thermocouples (colored by depth level: Front=red, Middle=blue,
  Back=green), 96~FCUs (gray squares), and ESS~113 (gold star).
  Arrows indicate the gradient-fit directions.
  Right: side cross-section with $z_\mathrm{norm}$ axis (0=back,
  1=front), showing the thinner M3 and thicker M1 regions.}
\label{fig:sensormap}
\end{figure}

%----------------------------------------------------------------------
\section{Surface--Air Coupling: Pre-Control}
\label{sec:pre_control}

Before the EAS activates the glycol setpoint, the mirror temperature
is governed solely by convective coupling with the cell air (ESS~113).
This pre-control period---typically several hours before sunset---provides
a clean measurement of the air--mirror thermal relationship at each
depth level.

\subsection{Pre-Control Offsets}

Figure~\ref{fig:pre_coupling} (top right) shows the distribution of
$T_\mathrm{surface} - \mathrm{ESS\,113}$ for each surface during
the pre-control period.  All three surfaces are close to ESS~113,
with small systematic offsets:
\begin{itemize}
  \item $T_\mathrm{front}$ is closest to ESS~113 (the optical surface
    is in direct convective contact with cell air).
  \item $T_\mathrm{back}$ tends to be slightly cooler (residual
    glycol cooling from the previous night).
\end{itemize}

The top-left panel shows $T_\mathrm{front}$ vs ESS~113 colored by
the through-glass gradient ($z\_gradient$), confirming tight coupling
between the optical surface and cell air when the glycol is off.
Per-night median offsets (bottom-left boxplot) show the inter-night
variability, and the bottom-right panel shows an example pre-control
night with all three surfaces tracking ESS~113 closely.

\begin{figure}[htbp]
\centering
\includegraphics[width=\textwidth]{../figures/mirror_surface_coupling_pre.png}
\caption{Surface--air coupling during pre-control (no glycol).
  Top left: $T_\mathrm{front}$ vs ESS~113 scatter colored by
  $z\_gradient$.  Top right: offset distributions for all three
  surfaces.  Bottom left: per-night median offset boxplots.
  Bottom right: example night showing $T_\mathrm{back}$,
  $T_\mathrm{mid}$, $T_\mathrm{front}$, and ESS~113.}
\label{fig:pre_coupling}
\end{figure}

%----------------------------------------------------------------------
\section{Surface--Air Coupling: Active Control}
\label{sec:active_control}

Once the glycol activates, the setpoint drives the back surface
away from air equilibrium, creating a through-glass gradient.
Figure~\ref{fig:active_coupling} characterizes this behavior.

\subsection{Surface Offset Evolution}

The top-right panel of Figure~\ref{fig:active_coupling} shows the
median (with IQR band) of $T_\mathrm{surface} - \mathrm{ESS\,113}$
as a function of time since setpoint activation.  Key observations:
\begin{itemize}
  \item All surfaces are driven below ESS~113 by the glycol cooling,
    consistent with the $-1.0$\,\textdegree C cold bias.
  \item $T_\mathrm{back}$ responds first and most strongly (closest
    to the FCUs).
  \item $T_\mathrm{front}$ lags and has a smaller offset, as the
    thermal gradient must propagate through the glass.
  \item The offsets stabilize after $\sim$3--4\,h, consistent with
    the $\tau \approx 3$\,h time constant.
\end{itemize}

\subsection{Through-Glass Gradient vs Setpoint Rate}

The bottom-left panel shows $z\_gradient$ vs the setpoint rate of
change.  When the setpoint is decreasing rapidly (aggressive cooling),
$z\_gradient$ tends negative---the back surface cools faster than the
front, as expected from the FCU proximity.  When the setpoint is
stable, $z\_gradient$ clusters near zero.

\begin{figure}[htbp]
\centering
\includegraphics[width=\textwidth]{../figures/mirror_surface_coupling_active.png}
\caption{Surface--air coupling during active control.  Top left:
  $T_\mathrm{front}$ vs ESS~113 colored by hours since activation.
  Top right: surface offset evolution (median $\pm$ IQR across nights).
  Bottom left: $z\_gradient$ vs setpoint rate (hexbin).
  Bottom right: example night with all surfaces and setpoint.}
\label{fig:active_coupling}
\end{figure}

%----------------------------------------------------------------------
\section{Glycol Response by Surface Depth}
\label{sec:glycol_depth}

The glycol system acts on the back surface through the FCUs.  The
thermal signal must propagate through the glass to reach the optical
surface.  We characterize this depth-dependent response using
per-surface cross-correlation lag and time constant fits.

\subsection{Cross-Correlation Lag}

Figure~\ref{fig:glycol_depth} (top left) shows the distribution of
cross-correlation lag between the setpoint rate and the rate of change
of each surface temperature.  The back surface responds first
(shortest lag), and the front surface last, as expected from the
thermal diffusion through the glass.

\subsection{Time Constant by Depth}

We fit the glycol-only first-order model
$dT_\mathrm{surface}/dt = (T_\mathrm{sp} - T_\mathrm{surface})/\tau$
independently for each surface level.  The results
(Figure~\ref{fig:glycol_depth}, top right) show:

\begin{table}[htbp]
\centering
\caption{Per-surface time constant from glycol-only first-order fit.}
\label{tab:tau_surface}
\begin{tabular}{lcc}
\toprule
Surface & $\tau$ median (h) & $\tau$ mean (h) \\
\midrule
$T_\mathrm{back}$  ($z_\mathrm{norm}=0$)   & 3.0 & 3.8 \\
$T_\mathrm{mid}$   ($z_\mathrm{norm}=0.5$) & 2.5 & 3.4 \\
$T_\mathrm{front}$ ($z_\mathrm{norm}=1.0$) & 2.0 & 2.9 \\
\bottomrule
\end{tabular}
\end{table}

The counterintuitive result that $\tau_\mathrm{front} < \tau_\mathrm{back}$
reflects that the front surface is also coupled to cell air convection
(ESS~113), which acts as a second thermal input.  The glycol-only model
absorbs this additional coupling into a shorter effective $\tau$.
The back surface, shielded from air convection by the glass itself,
is governed almost entirely by the glycol time constant.

\subsection{Glycol Forcing and Through-Glass Gradient}

The bottom panels of Figure~\ref{fig:glycol_depth} show the
relationship between glycol forcing and the through-glass gradient:
\begin{itemize}
  \item Bottom left: $z\_gradient$ vs (setpoint $-$ $T_\mathrm{back}$).
    When the glycol is actively cooling (setpoint $\ll T_\mathrm{back}$),
    the back cools faster than the front, creating a positive
    $z\_gradient$ (front warmer than back).
  \item Bottom right: $z\_gradient$ vs coolant supply--return
    $\Delta T$.  Strong cooling (large negative $\Delta T_\mathrm{coolant}$)
    correlates with positive $z\_gradient$.
\end{itemize}

\begin{figure}[htbp]
\centering
\includegraphics[width=\textwidth]{../figures/mirror_glycol_surface_response.png}
\caption{Glycol response by surface depth.  Top left: cross-correlation
  lag (setpoint rate $\to$ surface rate) per depth level.  Top right:
  $\tau_\mathrm{back}$ vs $\tau_\mathrm{front}$ scatter per night.
  Bottom left: $z\_gradient$ vs glycol forcing (setpoint $-$ $T_\mathrm{back}$).
  Bottom right: $z\_gradient$ vs coolant heat flux proxy
  (supply $-$ return $\Delta T$).}
\label{fig:glycol_depth}
\end{figure}

%----------------------------------------------------------------------
\section{Gradient Overview: Example Nights}
\label{sec:gradient_overview}

Figure~\ref{fig:gradient_overview} shows four example nights selected
by the magnitude of the through-glass gradient $|z\_gradient|$:
largest, smallest, median, and 25th percentile.  Each panel shows the
setpoint, ESS~113, and the three surface temperatures
($T_\mathrm{back}$, $T_\mathrm{mid}$, $T_\mathrm{front}$) with the
$\pm z\_gradient\_err$ uncertainty band shaded on $T_\mathrm{front}$.

Key observations:
\begin{itemize}
  \item On nights with large $|z\_gradient|$, the front and back
    surfaces separate by $\sim$0.3--0.5\,\textdegree C during active
    cooling, with the front remaining warmer.
  \item On nights with small $|z\_gradient|$, all three surfaces
    track each other closely, indicating near-isothermal conditions
    through the glass.
  \item The $z\_gradient\_err$ band is typically
    $\pm$0.1--0.2\,\textdegree C, confirming that the gradient
    decomposition is well-constrained by the 96 thermocouples.
\end{itemize}

\begin{figure}[htbp]
\centering
\includegraphics[width=\textwidth]{../figures/mirror_gradient_overview.png}
\caption{Gradient overview: four example nights selected by
  $|z\_gradient|$ magnitude.  Each panel shows setpoint (black),
  ESS~113 (dotted), and $T_\mathrm{back}$ (green), $T_\mathrm{mid}$
  (blue), $T_\mathrm{front}$ (red) with $\pm z\_gradient\_err$
  shaded band.}
\label{fig:gradient_overview}
\end{figure}

%----------------------------------------------------------------------
\section{Setpoint--Surface Coupling: Inner Ring vs Outer Ring}
\label{sec:ring_coupling}

The gradient-fit decomposition (Section~\ref{sec:glycol_depth}) captures
the average through-glass behavior, but the 96~thermocouples span a
wide range of radii---from the thin M3 inner region ($r \approx 0.56$\,m,
glass thickness 0.24\,m) to the thick M1 outer region
($r > 3.5$\,m, glass thickness 0.74\,m).  To resolve the radial
dependence of the glycol coupling, we use \textbf{raw thermocouple
readings} (not the gradient-fit decomposition) averaged over two
concentric rings at two depth levels.

\subsection{Data and Method}

The inner ring uses 14 MTCIW sensors at $r \approx 0.56$\,m (6~front,
2~middle, 6~back), averaged into \texttt{inner\_front} and
\texttt{inner\_back}.  The outer ring uses all sensors at $r > 3.5$\,m,
similarly averaged.  These are merged with the ESS~113, setpoint, and
dome status from the thermal parquet (July 2025+) and a supplementary
ESS~113 query (May--June 2025), yielding 155 valid nights with
concurrent setpoint and thermocouple data.

For each night and each of the four surfaces (inner front, inner back,
outer front, outer back), we:
\begin{enumerate}
  \item Fit the first-order model $dT/dt = (T_\mathrm{sp} - T)/\tau$
    by grid search over $\tau \in [0.5, 12]$\,h in 0.25\,h steps.
  \item Compute the cross-correlation lag between the setpoint rate
    and the surface temperature rate.
\end{enumerate}

\subsection{Time Constant by Ring and Depth}

\begin{table}[htbp]
\centering
\caption{Setpoint-to-surface thermal coupling: raw thermocouple results
  (155~nights, cleaned setpoint---gaps $> 5$\,min excluded).}
\label{tab:ring_coupling}
\begin{tabular}{lccc}
\toprule
Surface & $\tau$ median (h) & Lag median (min) & RMSE median (\textdegree C) \\
\midrule
Inner front ($r \approx 0.56$\,m, $z=1$) & 2.00 & 4 & 0.098 \\
Inner back  ($r \approx 0.56$\,m, $z=0$) & 3.00 & 9 & 0.066 \\
Outer front ($r > 3.5$\,m, $z=1$) & 2.00 & 3 & 0.105 \\
Outer back  ($r > 3.5$\,m, $z=0$) & 3.00 & 6 & 0.070 \\
\bottomrule
\end{tabular}
\end{table}

The results (Table~\ref{tab:ring_coupling}, Figure~\ref{fig:ring_tau})
reveal clear physical trends:
\begin{itemize}
  \item \textbf{Front surfaces respond faster than back surfaces} at
    both rings (inner: 2.0 vs 3.0\,h; outer: 2.0 vs 3.0\,h).
    This confirms the gradient-fit finding (Section~\ref{sec:glycol_depth})
    and reflects the dual coupling of the front surface to both glycol
    and cell air convection.
  \item \textbf{Back surfaces converge at $\tau = 3.0$\,h} regardless
    of ring radius, suggesting that the glycol-to-back coupling is
    dominated by the FCU contact resistance rather than glass thickness.
  \item \textbf{The back surfaces are the thermal bottleneck}:
    $\tau = 3.0$\,h at both rings, with the outer back having the
    highest RMSE (0.070\,\textdegree C) and longest lag (6\,min).
  \item \textbf{Front surfaces respond in $\sim$2\,h} at both rings,
    with lags of only 3--4\,min, due to the additional convective
    coupling with cell air.
  \item \textbf{Inner back has the lowest RMSE} (0.066\,\textdegree C),
    confirming that the thin-glass back surface most directly reflects
    the glycol setpoint.
\end{itemize}

\begin{figure}[htbp]
\centering
\includegraphics[width=\textwidth]{../figures/setpoint_surface_tau_lag.png}
\caption{Setpoint-to-surface coupling: $\tau$ and lag.  Top left:
  $\tau$ distributions by surface.  Top right: $\tau$ boxplots by
  ring $\times$ depth.  Bottom left: cross-correlation lag distributions.
  Bottom right: $\tau_\mathrm{inner\_back}$ vs
  $\tau_\mathrm{outer\_back}$ (same depth, different ring thickness).}
\label{fig:ring_tau}
\end{figure}

\subsection{Through-Glass Gradient by Ring}

The through-glass temperature difference (front $-$ back) is near zero
on average at both rings (inner median: $+0.002$\,\textdegree C, outer
median: $-0.006$\,\textdegree C), but the inner ring shows larger
scatter ($\sigma = 0.18$\,\textdegree C vs $0.15$\,\textdegree C for
the outer ring).

Figure~\ref{fig:ring_gradient} (top right, bottom left) shows that
when the glycol is strongly cooling (setpoint $\ll T_\mathrm{back}$),
the inner ring develops a clear negative front-back gradient (back
cooling faster than front), reaching $\pm 0.5$\,\textdegree C.
The outer ring shows a tighter gradient range ($\pm 0.3$\,\textdegree C),
as the thicker glass dampens the through-thickness temperature difference.

\begin{figure}[htbp]
\centering
\includegraphics[width=\textwidth]{../figures/setpoint_glass_gradient.png}
\caption{Through-glass gradient: inner ring (thin, 0.24\,m) vs outer
  ring (thick, 0.74\,m).  Top left: front$-$back $\Delta T$
  distributions.  Top right: inner ring glass $\Delta T$ vs glycol
  forcing (hexbin).  Bottom left: outer ring glass $\Delta T$ vs glycol
  forcing.  Bottom right: per-night median inner vs outer gradient.}
\label{fig:ring_gradient}
\end{figure}

\subsection{Gradient--Forcing Slope: Safe Setpoint Range}

To determine the glycol forcing that can be applied without creating
excessive through-glass gradients, we fit the relationship
\begin{equation}
  \Delta T_\mathrm{glass} = T_\mathrm{front} - T_\mathrm{back}
    = \alpha \cdot (T_\mathrm{sp} - T_\mathrm{back}) + \beta
  \label{eq:grad_slope}
\end{equation}
using Theil-Sen robust regression (resistant to outliers) on 80{,}361
minute-level samples per ring (Figure~\ref{fig:grad_slope}).

\begin{table}[htbp]
\centering
\caption{Theil-Sen robust fit: through-glass gradient vs glycol forcing.}
\label{tab:grad_slope}
\begin{tabular}{lccc}
\toprule
Gradient measure & Slope $\alpha$ (\textdegree C/\textdegree C) & Intercept $\beta$ (\textdegree C) & $N$ samples \\
\midrule
$z\_gradient$ (96-TC fit) & 0.223 & $-0.011$ & 86{,}472 \\
Inner ring (0.24\,m glass) & 0.212 & $+0.006$ & 80{,}361 \\
Outer ring (0.74\,m glass) & 0.181 & $+0.004$ & 80{,}361 \\
\bottomrule
\end{tabular}
\end{table}

The slope $\alpha$ gives the fraction of the glycol forcing that
appears as a through-glass gradient.  All three measures are consistent:
$\alpha \approx 0.18$--$0.22$\,\textdegree C/\textdegree C.  The
96-thermocouple $z\_gradient$ (which averages over all radii) has the
steepest slope (0.223), while the inner ring raw TCs give 0.212 and
the outer ring gives 0.181---the thinner inner glass transmits more of the
forcing as a gradient.

Using these slopes, we can infer the safe forcing range for a given
gradient threshold:

\begin{table}[htbp]
\centering
\caption{Maximum glycol forcing ($T_\mathrm{sp} - T_\mathrm{back}$)
  to stay within gradient limits.}
\label{tab:safe_forcing}
\begin{tabular}{lccc}
\toprule
$|\Delta T_\mathrm{glass}|$ limit & $z\_gradient$ forcing & Inner ring forcing & Outer ring forcing \\
\midrule
0.1\,\textdegree C & $\pm$0.4\,\textdegree C & $\pm$0.5\,\textdegree C & $\pm$0.6\,\textdegree C \\
0.2\,\textdegree C & $\pm$0.8\,\textdegree C & $\pm$1.0\,\textdegree C & $\pm$1.1\,\textdegree C \\
0.3\,\textdegree C & $\pm$1.3\,\textdegree C & $\pm$1.4\,\textdegree C & $\pm$1.7\,\textdegree C \\
\bottomrule
\end{tabular}
\end{table}

For the typical EAS cold bias of $-1.0$\,\textdegree C, the expected
glass gradient is $\sim$0.2\,\textdegree C (all measures)---well within
safe limits.  Aggressive transient forcing of $-5$\,\textdegree C (as
seen during rapid cool-downs) would produce $\sim$1\,\textdegree C
gradients at the inner ring.

\begin{figure}[htbp]
\centering
\includegraphics[width=\textwidth]{../figures/setpoint_gradient_robust_fit.png}
\caption{Through-glass gradient vs glycol forcing
  ($T_\mathrm{sp} - T_\mathrm{back}$) with Theil-Sen robust linear
  fit (red line).  Left: 96-TC gradient-fit $z\_gradient$
  (slope = 0.223).  Center: inner ring raw TCs (0.24\,m glass,
  slope = 0.212).  Right: outer ring raw TCs (0.74\,m glass,
  slope = 0.181).}
\label{fig:grad_slope}
\end{figure}

\subsection{Glycol Coolant: Active vs Passive Operation}
\label{sec:coolant_coupling}

The glycol coolant loop has 97\% time coverage, while setpoint commands
cover only 8.9\% of the data.  The coolant circulates continuously,
but does its heat flux actually drive the mirror temperature?

\subsubsection{Coolant $\Delta T$ Does Not Predict Mirror Cooling Rate}

We test the causal relationship by correlating the coolant heat flux
proxy ($\Delta T_\mathrm{coolant} = T_\mathrm{supply} - T_\mathrm{return}$)
with the 5-minute smoothed mirror cooling rate
$dT_\mathrm{mirror}/dt$ (Figure~\ref{fig:coolant}).

\begin{table}[htbp]
\centering
\caption{Pearson correlation $r$ between mirror cooling rate
  ($dT/dt$ inner back) and candidate drivers.}
\label{tab:coolant_corr}
\begin{tabular}{lcc}
\toprule
Driver & Period & $r$ \\
\midrule
Setpoint $-$ $T_\mathrm{back}$ & Active & \textbf{0.577} \\
Coolant $\Delta T$ (supply $-$ return) & Active & 0.075 \\
Coolant $\Delta T$ (supply $-$ return) & Passive & 0.024 \\
\bottomrule
\end{tabular}
\end{table}

The results (Table~\ref{tab:coolant_corr}) are unambiguous:
\begin{itemize}
  \item The \textbf{setpoint error strongly predicts} the mirror cooling
    rate ($r = 0.577$, Figure~\ref{fig:coolant} bottom left).  When
    the setpoint is far below the mirror, the mirror cools at up to
    $-1.5$\,\textdegree C/h; near equilibrium, $dT/dt \approx 0$.
    The Theil-Sen slope is 0.24\,(°C/h)/°C, consistent with
    $\tau \approx 1/0.24 \approx 4$\,h.
  \item The \textbf{coolant $\Delta T$ does not predict} the mirror cooling
    rate---neither during active ($r = 0.075$) nor passive
    ($r = 0.024$) operation.  Despite coolant $\Delta T$ reaching
    $-6$\,\textdegree C during active cooling, it shows essentially
    no linear correlation with $dT_\mathrm{mirror}/dt$.
\end{itemize}

\subsubsection{Interpretation}

The coolant $\Delta T$ reflects how much heat the glycol \emph{has
already absorbed} from the mirror, not how much cooling it will
\emph{deliver next}.  The 96~FCUs modulate their local flow and fan
speed based on the commanded setpoint---the setpoint governs
\emph{how much} of the available chiller capacity is directed at
the mirror.  The coolant supply--return $\Delta T$ is therefore a
\emph{consequence} of the setpoint-driven cooling, not its cause.

\subsubsection{Energy Balance}

Despite the lack of predictive correlation, the energy balance
relates the coolant $\Delta T$ to the mirror's thermal state.
The heat extraction rate through the glycol loop is:
\begin{equation}
  \dot{Q} = \dot{m}\, c_p\, (T_\mathrm{return} - T_\mathrm{supply})
  \label{eq:qdot}
\end{equation}
where $\dot{m}$ is the mass flow rate and $c_p$ the specific heat.
For 50\% propylene glycol at $\sim$10\,\textdegree C:
$\rho = 1040$\,kg/m$^3$, $c_p = 3560$\,J/(kg\,K).

Equating coolant heat extraction to the mirror's temperature change:
\begin{equation}
  \dot{m}\, c_p\, |\Delta T_\mathrm{coolant}|
    = M_\mathrm{mirror}\, c_\mathrm{glass}\,
      \left|\frac{dT_\mathrm{mirror}}{dt}\right|
  \label{eq:mirror_energy}
\end{equation}
For M1M3: $M_\mathrm{mirror} \approx 30{,}000$\,kg,
$c_\mathrm{glass} \approx 670$\,J/(kg\,K) (borosilicate).
From $\tau = 3$\,h, the effective thermal conductance is:
\begin{equation}
  h_\mathrm{eff}\, A
    = \frac{M_\mathrm{mirror}\, c_\mathrm{glass}}{\tau}
    = \frac{30{,}000 \times 670}{10{,}800}
    \approx 1860\;\mathrm{W/K}
  \label{eq:heff}
\end{equation}
At a typical cooling rate of $-1.0$\,\textdegree C/h, the required
heat extraction is $\dot{Q} \approx 5.6$\,kW.  With
$\Delta T_\mathrm{coolant} = -4$\,\textdegree C, this implies
$\dot{m} \approx 0.4$\,kg/s ($\sim$6\,GPM across 96~FCUs).

\subsubsection{Passive Circulation}

During passive operation (no setpoint, 91\% of data),
$\Delta T_\mathrm{coolant} \approx 0$\,\textdegree C
(Figure~\ref{fig:coolant}, top left).  The glycol circulates to
maintain hydraulic readiness but the FCUs are throttled to minimum
flow, extracting negligible heat.  The mirror temperature during
these periods drifts primarily in response to ambient air convection,
not the coolant system ($r = 0.024$).

\begin{figure}[htbp]
\centering
\includegraphics[width=\textwidth]{../figures/coolant_surface_coupling.png}
\caption{Does coolant heat flux drive the mirror?  Top left: coolant
  $\Delta T$ distribution---passive (green, near zero) vs active (blue,
  reaching $-6$\,\textdegree C).  Top right: $dT/dt$ vs coolant
  $\Delta T$ during active control ($r = 0.075$, no correlation).
  Bottom left: $dT/dt$ vs setpoint error ($r = 0.577$, strong
  correlation).  Bottom right: $dT/dt$ vs coolant $\Delta T$ during
  passive circulation ($r = 0.024$, no correlation).
  The setpoint error is the driver; the coolant $\Delta T$ is a
  consequence.}
\label{fig:coolant}
\end{figure}

\subsection{Example Nights and Model Quality}

Figure~\ref{fig:ring_examples} shows four example nights: fastest and
slowest outer-ring $\tau$, best and median model fit.  The first-order
model (dashed black) captures the thermal response well for most nights.
On nights with large $\tau$ (top right), the surfaces barely respond to
the setpoint---these are typically nights where the setpoint changes
are small and the mirror is already near equilibrium.

Figure~\ref{fig:ring_rmse} summarizes the model fit quality.  The
RMSE distributions (top left) cluster below 0.25\,\textdegree C for
all four surfaces, confirming that the first-order model captures
the dominant physics even at individual thermocouple locations.  The
$\tau_\mathrm{inner\_front}$ vs $\tau_\mathrm{outer\_front}$ scatter
(bottom left) shows that the outer optical surface consistently
requires a longer time constant, reflecting the thicker glass path.

\begin{figure}[htbp]
\centering
\includegraphics[width=\textwidth]{../figures/setpoint_surface_examples.png}
\caption{Four example nights showing setpoint (black), ESS~113 (dotted),
  and all four surfaces (inner front/back, outer front/back).  The
  dashed line shows the first-order model fit for the outer back
  surface.  Panels selected by: fastest $\tau$ (top left), slowest
  $\tau$ (top right), best RMSE (bottom left), median RMSE
  (bottom right).}
\label{fig:ring_examples}
\end{figure}

\begin{figure}[htbp]
\centering
\includegraphics[width=\textwidth]{../figures/setpoint_surface_rmse.png}
\caption{Model fit quality by ring $\times$ depth.  Top left: RMSE
  distributions.  Top right: RMSE boxplots.  Bottom left:
  $\tau_\mathrm{inner\_front}$ vs $\tau_\mathrm{outer\_front}$ scatter.
  Bottom right: RMSE inner front vs RMSE outer front.}
\label{fig:ring_rmse}
\end{figure}

%----------------------------------------------------------------------
\section{Implications for Control Simulation}

\begin{enumerate}
  \item \textbf{Cold bias should be 1.0\,\textdegree C}, not
    0.3\,\textdegree C.  The EAS targets $\mathrm{ESS\,113} - 1.0$\,\textdegree C.

  \item \textbf{Bulk time constant is $\tau = 2.5$\,h} (for $T_\mathrm{mid}$).
    Surface-dependent: $\tau_\mathrm{back} = 3.0$\,h,
    $\tau_\mathrm{front} = 2.0$\,h.

  \item \textbf{Rate limit of 0.7\,\textdegree C/h is conservative.}
    The EAS routinely exceeds this (p95 = 1.8\,\textdegree C/h),
    confirming that the soft-rate-penalty approach (allowing occasional
    violations up to $\sim$2\,\textdegree C/h) is physically realistic.

  \item \textbf{The first-order model is excellent for the bulk mirror}
    when validated against $T_\mathrm{mid}$ ($z_\mathrm{norm} = 0.5$).

  \item \textbf{Current EAS achieves 26.5\% within $\pm$0.3\,\textdegree C.}
    This is the baseline to beat with improved forecasting and control.

  \item \textbf{The optical surface ($T_\mathrm{front}$) is coupled to
    both glycol and cell air.}  Its effective $\tau_\mathrm{front} = 2.0$\,h
    is shorter than $\tau_\mathrm{back} = 3.0$\,h because convective
    coupling with ESS~113 provides a second thermal input.  A two-input
    model may better represent the front-surface dynamics.

  \item \textbf{The through-glass gradient $z\_gradient$ is a diagnostic
    for control aggressiveness.}  Large $|z\_gradient|$ ($> 0.3$\,\textdegree C)
    indicates the glycol is driving the back surface faster than the
    front can follow.  Future controllers could monitor $z\_gradient$
    as a constraint to limit thermal stress.

  \item \textbf{Back surfaces are the thermal bottleneck}
    ($\tau = 3.0$\,h at both rings), while front surfaces respond in
    $\sim$2\,h due to additional air convection coupling.  The inner
    back has the lowest RMSE (0.070\,\textdegree C), confirming it
    most directly reflects the glycol setpoint.

  \item \textbf{Glass thickness affects the front surface more than the
    back.}  Back-surface $\tau$ is the same at both rings (3.0\,h),
    suggesting FCU contact resistance dominates over glass conduction
    for the glycol-facing surface.
\end{enumerate}

%----------------------------------------------------------------------
\section{Conclusions}

\begin{enumerate}
  \item The EAS targets $T_\mathrm{setpoint} = \mathrm{ESS\,113} -
    1.0$\,\textdegree C, updating every 3\,min at rates up to
    1.8\,\textdegree C/h (p95).

  \item The glycol system responds to setpoint changes with a
    4-minute lag (at $T_\mathrm{mid}$).  The bulk glass temperature
    tracks the setpoint within $-0.07$\,\textdegree C (median offset).

  \item The first-order thermal model with $\tau = 2.5$\,h
    achieves median RMSE = 0.059\,\textdegree C against
    $T_\mathrm{mid}$---an excellent fit that validates the model
    structure.

  \item The bulk glass time constant is $\tau \approx 2.5$\,h
    ($T_\mathrm{mid}$), with depth-dependent values from
    $\tau_\mathrm{back} = 3.0$\,h to $\tau_\mathrm{front} = 2.0$\,h.

  \item The current EAS achieves per-night RMS = 0.74\,\textdegree C
    and 26.5\% within $\pm$0.3\,\textdegree C---significant room for
    improvement with better forecasting and control algorithms.

  \item The gradient-fit decomposition resolves three depth levels
    through the glass.  The effective time constant decreases from
    $\tau_\mathrm{back} = 3.0$\,h to $\tau_\mathrm{front} = 2.0$\,h,
    reflecting the dual coupling of the optical surface to both glycol
    and cell air convection.

  \item The through-glass gradient $z\_gradient$ is typically
    $< 0.3$\,\textdegree C but can reach 0.5\,\textdegree C during
    aggressive cooling.  The gradient uncertainty
    ($z\_gradient\_err \approx 0.1$--0.2\,\textdegree C) confirms
    the decomposition is well-constrained.

  \item Pre-control (no glycol), all three surfaces track ESS~113
    closely with small offsets, confirming good air--mirror thermal
    equilibrium before sunset.

  \item Raw thermocouple analysis of inner ring ($r \approx 0.56$\,m,
    glass thickness 0.24\,m) vs outer ring ($r > 3.5$\,m, 0.74\,m)
    over 155~nights (with cleaned setpoint signal) shows: back
    surfaces converge at $\tau = 3.0$\,h regardless of radius (FCU
    contact resistance dominates), while front surfaces respond in
    $\sim$2\,h due to air convection coupling.  Model RMSE is
    0.07--0.10\,\textdegree C across all four ring$\times$depth
    combinations.

  \item Theil-Sen robust regression of the through-glass gradient
    vs glycol forcing yields slopes of 0.21\,\textdegree C/\textdegree C
    (inner) and 0.18\,\textdegree C/\textdegree C (outer).  This
    means the EAS's typical $-1$\,\textdegree C forcing produces
    only $\sim$0.2\,\textdegree C glass gradient.  To keep gradients
    below 0.3\,\textdegree C, forcing should stay within
    $\pm$1.4\,\textdegree C (inner) or $\pm$1.7\,\textdegree C (outer).
\end{enumerate}

\end{document}
